\documentclass{article}
\usepackage[a4paper, margin=1in]{geometry}
\usepackage{graphicx} 
\usepackage{amsmath}
\usepackage{amssymb}
\usepackage{enumitem}

\title{\vspace{-2cm} \textbf{POE Question Bank \\ Semester II 2024-2025}}
\begin{document}
\date{}
\maketitle
\vspace{-1.5cm}
\noindent\hrulefill
\section*{Part A}
\begin{enumerate}
    \item 
        Quantity demanded for a good is given by the function
        \[
        Q_A^d = 40 - 0.02P_A + 0.9P_B
        \]
        Where:
        \begin{itemize}
            \item $P_A$ = price of commodity $A$
            \item $P_B$ = price of commodity $B$
        \end{itemize}
        What is the quantity demanded when the prices of good $A$ and $B$ are \$20 and \$8 respectively?        
    \item 
        A consumer's Utility Function is given by
        \[
        U(x,y) = x^2y
        \]
        What is the total utility obtained from a consumption bundle of (3,2)?
    \item 
        The price of a product increases from \$5 to \$7, and the quantity demanded decreases from 100 units to 80 units. Calculate the price elasticity of demand. 
    \item 
        A producer is willing to sell a good for \$15 but manages to sell it at the market price of \$25. Calculate the producer surplus per unit.
    \item 
        Cost Function for a firm is given by 
        \[
        C = 5Q^3 + 8Q - ln(Q)
        \]
        What is the cost incurred by firm to produce one additional unit at a production level of 12 units?
    \item
        A firm producing commodity A incurs fixed costs of \$40 and average variable costs of \$5. Market price of A is \$10. How many units must the firm sell to break even?
    \item 
        Demand and Supply Functions are given by 
        \begin{align}
        Q^d = 16-0.2P\\
        Q^s = 0.4P-11
        \end{align}
        Find the equilibrium price and quantity.
    \item 
        A firm produces in perfectly competitive market where the price is \$10. If the firm's cost function is given by
        \[
        C=(Q+3)^2
        \]
        Find the profit maximizing quantity.
    \item
        The supply function of a firm is 
        \[
        Q^s = 2P-10
        \]
        If the market clearing price is \$20, calculate the equilibrium quantity and producer surplus. 
    \item 
        A firm produces two goods, X and Y. The cross-price elasticity of demand between them is -0.8. Comment on the nature of these goods. If the price of X increases by 10\%, what is the percentage change in the quantity demanded of Y?
\end{enumerate} 

\section*{Part B}
\begin{enumerate}
    \item
        If the demand function for good $x$ is given by
        \[
        Q^d = 6 - 2P_x + 4P_y
        \]
        where the prices of good $y$ are fixed at \$1.
        Find the inverse demand function and consumer surplus if the price of good $x$ is \$3?
    \item 
        Quantity demanded for a good is given by the function
        \[
        Q_A^d = 40 - 0.02P_A + 0.9P_B + 0.3Y_d
        \]
        Where:
        \begin{itemize}
            \item $P_A$ = price of commodity $A$
            \item $P_B$ = price of commodity $B$
            \item $Y_d$ = disposable income 
        \end{itemize}
        Find price elasticity of demand when the price of good $A$ is \$20, price of good $B$ is \$8 and consumer's disposable income is \$500.
    \item 
        A firm incurs fixed costs of \$300. Variable Cost function is given by 
        \[
        C_v = 50Q^2 + 13Q
        \]
        and the Total Revenue function is given by
        \[
        R = 10Q^3 - 4Q^2
        \]
        The firm produces $6$ units. Calculate firm's loss and average total cost. Should the firm continue to operate in short run?
    \item 
        The price elasticity of demand and price elasticity of supply for a good are $-0.4$ and $0.5$ respectively at an equilibrium price of \$5 and quantity of 16 units. Find the demand and supply functions.  
    \item 
        Utility obtained from consumptions of good X and Y can be expressed as
        \[
        U(x,y) = xy
        \]
        If the consumer has \$40, and the prices of X and Y are \$4 and \$2, respectively, what is their optimum consumption bundle?
    \item 
        Bob quits his computer programming job, where he was earning a salary of \$120,000 per year, to start his own business in a building he owns and was previously renting out for \$30,000 per year. He incurs the following expenses
        \begin{itemize}
            \item Salary paid to himself - \$55,000
            \item Other Expenses - \$25,000
        \end{itemize}
        Find the accounting cost and economic cost associated with his business.
    \item 
        A company using only Labor(L) and Capital(K) to manufacture its goods follows the Cobb-Douglas production function given by 
        \[
        Q = 6L^\frac{1}{2}K^\frac{2}{3}
        \]
        The company projects the total cost of production to be \$4200. Per unit cost of labor and capital are \$20 and \$120 respectively. What combination of capital and labor should it use?
    \item 
        A firm's production function is given by 
        \[
        F(K,L) = AK^\alpha L^\beta 
        \]
        where K is capital, L is labor, and $A,\alpha,\beta > 0$. Derive the conditions for which this function exhibits increasing, decreasing and constant returns to scale.
    \item 
        A monopolist faces the following demand curve
        \[
        P=100-2Q
        \]
        The monopolist's total cost function is given by 
        \[
        TC=20+3Q^2
        \]
        Find the profit maximizing price and quantity.
    \item 
        Demand and Supply Functions are given by 
        \begin{align}
        Q^d = 2-\frac{1}{4}P \tag{1}\\
        Q^s = \frac{1}{2}P - \frac{5}{2} \tag{2}
        \end{align}
        If government intervention results in a price floor at \$7, then calculate consumer surplus, producer surplus and deadweight loss.
\end{enumerate}

\section*{Part C}
\begin{enumerate}
    \item 
        Mike consumes two goods A and B. The utility he gains from consumption of these products is given by 
        \[
        U(A,B) = 4A^2B
        \]
        His monthly income is \$120 and prices of A and B are \$5 and \$8 respectively. Suppose the price of good A decreases by \$2, then calculate the following. 
        \begin{enumerate}[label=(\roman*)]
            \item Income Effect on A
            \item Substitution Effect on B
        \end{enumerate}        
    \item 
        Suppose Bob Auto Inc. is the only firm producing solar powered cars. It's demand and cost curves are as follows 
        \begin{align}
            P=30-3Q \tag{1} \\
            TC=3Q^2+5 \tag{2} 
        \end{align}
        Find the socially optimal level of production and the social cost of monopoly. 
    \item 
        In a perfectly competitive market, the demand and supply functions are
        \begin{align}
            Q^d = 200-5P \tag{1} \\
            Q^s = 3P-60 \tag{2} 
        \end{align}
        Suppose government imposes a \$10 per unit tax on the producers, then calculate 
        \begin{enumerate}[label=(\roman*)]
            \item Tax burden on Consumers
            \item Tax burden on Producers
        \end{enumerate}
    \item 
        The demand curve for a manufacturing company is given by 
        \[
        Q^d = 100-5P
        \]
        It  has two plants and the marginal cost for each plant is given by 
        \begin{align*}
        MC_1 = 10+5Q_1 \\
        MC_2 = 15+5Q_2
        \end{align*}
        Determine the profit-maximizing output level of each plant and equilibrium price. 
    \item 
        Suppose two goods X and Y are perfect complements of each other, i.e. they are consumed together in a fixed ratio. It is known that the ratio of consumption $x:y$ is $3:2$. Find U(x,y) and plot the indifference curve U(x,y) = 10.
    \item 
        Assume Pooja Wines is the only supplier in the local market. They face the following demand curve
        \[
        P=180-2Q
        \]
        Their total cost function is given by 
        \[
        TC=Q^2+20
        \]
        If the market structure was initially perfectly competitive with $P=\$90$, calculate the change in producer surplus, consumer surplus and deadweight loss.
    \item 
        If a firm is given a subsidy of \$5/unit produced, how would this affect the good's price? Also determine the benefit distribution between consumers and producers. It is known that the price elasticity of demand ($\epsilon_d$) and price elasticity of supply ($\epsilon_s$) is -0.8 and 1.2 respectively. 
    \item 
        A consumer's utility function is given by 
        \[
        U(x,y) = ln(x)+y
        \]
        This type of a utility function is also called the quasi-linear utility function. For an arbitrary budget line $M=p_xx+p_yy$, find an expression for the optimal consumption bundle $(x^*,y^*)$. Also plot and interpret the graph of $M$ against $x^*$.
    \item 
        In a perfectly competitive market, the demand and supply functions are
        \begin{align}
            Q^d = 200-5P \tag{1} \\
            Q^s = 3P-60 \tag{2} 
        \end{align}
        Suppose government imposes a \$10 per unit tax on the producers, then calculate 
        \begin{enumerate}[label=(\roman*)]
            \item Deadweight Loss from Taxation
            \item Tax Revenue
        \end{enumerate}
    \item 
        Consider the case for commodities X and Y which are perfect substitutes of each other, with a linear utility function given by
        \[
        U(x,y) = 3x+4y 
        \]
        The following information is known:
        \begin{itemize}
            \item Price of $X$ = \$2
            \item Price of $Y$ = \$5
            \item Income = \$10
        \end{itemize}
        Graphically interpret and find the equation of the indifference curve with maximum utility.
    \item 
        Suppose a consumer's preferences are represented by a utility function $U(x,y) = x + ln(y)$. The consumer has an income $I=50$, and initial prices are $P_x = 1$ and $P_y = 1$. The price of x increases from $P_x =1$ to $P_x = 1.5$. What type of a good is x?
        \begin{enumerate}[label=(\alph*)] 
            \item Giffen Good
            \item Normal Good
            \item Inferior Good
            \item None of the above
        \end{enumerate}
    \item 
        If a consumer's initial basket given by (x,y) is (5,7). Suppose the price of good x decreases and his new basket is (8,7). It is known that during the shift, his decomposition basket was (10,5). What type of a good is x?
        \begin{enumerate}[label=(\alph*)] 
            \item Giffen Good
            \item Normal Good
            \item Inferior Good
            \item None of the above
        \end{enumerate}
    \item 
        A consumer maximizes their utility $U(x,y) = x^{0.3}y^{0.7}$, with an income $I=120$ and initial prices are $P_x = 6$ and $P_y=4$. The price of x decreases from $P_x = 6$ to $P_x = 3$. What is x?
        \begin{enumerate}[label=(\alph*)] 
            \item Giffen Good
            \item Normal Good
            \item Inferior Good
            \item None of the above
        \end{enumerate}
    \item 
        Suppose Checks\kern0.1em\&Co.\kern0.4em is the only screw manufacturer in the market. Their inverse demand function is given by $P=20-Q$ and their marginal cost is given by $MC=2Q$. If the government imposes a \$4/unit tax on the monopolist, then calculate changes in 
        \begin{enumerate}[label=(\roman*)]
            \item Consumer Surplus
            \item Producer Surplus
            \item Deadweight Loss
        \end{enumerate}
    \item 
        The price elasticity of demand and price elasticity of supply for a good are $-0.4$ and $0.5$ respectively at an equilibrium price of \$5 and quantity of 16 units. Determine the new equilibrium price and quantity if demand curve shifts to the right by 10\%.
    \item 
        Consider the utility function
        \[
        U(x,y) = \sqrt{x^2+y^2}
        \]
        The price of good x is \$3, price of good y is \$4 and your monthly income is \$50.
        \begin{enumerate}[label=(\roman*)]
            \item Using the tangency condition, find the optimal consumption bundle.  
            \item Explain whether the bundle above truly maximizes utility. If not, what is the actual optimal consumption bundle?
        \end{enumerate} 
    \item 
        Given the utility function and budget constraint
        \begin{align*}
        U(x_1,x_2) = x_1^\alpha x_2^\beta \\
        p_1x_1+p_2x_2 = M
        \end{align*}
        Find an expression for $U(x_1^*,x_2^*)$ where $(x_1^*,x_2^*)$ represents the optimal consumption bundle. 
    \item 
        A firm's production function is given by
        \[
        Q=4K^\frac{2}{7}L^\frac{5}{7}
        \]
        and the input prices are r=\$12 and w=\$5. Determine the cost-minimizing combination of K and L if the firm wants to produce  150 units of output.
    \item 
        Adidas has just launched a new sneakers collection in the market. Due to uncertainty in the market, prices could either be $P_1=450-0.58Q$ or $P_2=350-0.85Q$. If probability of $P_1$ is 0.45, determine the optimal output to maximize expected profit if their marginal cost is constant and equal to 80.
    \item 
        A non linear demand function is given by $Q=aP^{-b}$, where a and b are positive constants. How does the point price elasticity of demand vary along the demand curve? If $\Delta P =-10\%$, find the percentage change in quantity demanded.
\end{enumerate}

\end{document}
